\scanpage{28}
\uline{Corollaire 6} $P/X/S$, avec $X\to S$ propre plat à fibres géométriques intègres. Supposons qu'il existe une section $\sigma$ de $X/S$ le long de laquelle $X/S$ soit \deleted{simple} \inserted{lisse} \deleted{et $P/X$ \unsure{est} \unsure{l.f.f.} \unsure{séparé}}. Alors $\prod_{X/S}P/X$ représentable. Si $P\mathbin{\times}_X(S,\sigma)$ est quasi-projectif sur $S$, alors $\prod_{X/S}P/X$ est essentiellement quasi-projectif.

\uline{Corollaire 7} $P/X/S$, avec $X\to S$ simple propre \deleted{à fibres géométriques intègres} \inserted{connexes, $P/X$ \unsure{l.f.f.|l.f.t.} séparé}. Alors $\prod_{P/X}X/S$ est localement représentable pour la top. étale.

Remplaçant la top étale par la top. \unsure{fpqf} (fid plt quasi fini), le résultat 6 demeure être valable sans supposer $X\to S$ simple, mais seulement propre, plat, à fibres \deleted{géométriques} \inserted{sans cycles} associés immergés, par essentiellement la $=$ démonstration.

\uline{Corollaire 8} $X$ et $Y$ propres sur corps $k$. Alors $\unsure{\underline{\operatorname{Isom}}}_k(X,Y)$ est représentable. En particulier, $\underline{\operatorname{Aut}}_k(X)$ est représentable.
% archive foliation: 27
